
\needspace{4\baselineskip}
\color{uniblue}{\subsection{Контроль токовых цепей (КЦТ)}\label{sec:kzt}}
\color{black}

\begin{enumerate}[label=\arabic{section}.\arabic{subsection}.\arabic*, labelsep=4pt, leftmargin=0pt, itemindent=57pt]

\item
КЦТ предназначен для контроля токовых цепей, подключенных к аналоговым входам устройства.

\item
В устройстве контроль токовых цепей представлен функциональным блоком «КЦТ». В состав функционального блока «КЦТ» входят функции контроля токовых цепей по сторонам трансформатора (КЦТ ВН, КЦТ НН1, КЦТ НН2). Сигнал неисправности токовых цепей формируется логическим суммированием сигналов нарушения симметрии от каждой из функций контроля токовых цепей сторон трансформатора. При фиксации неисправности токовых цепей формируется сигнал «КЦТ: Срабатывание», который заводится в ДТЗт (см. пункт \ref{sec:dtzt}). 

\item
КЦТ может быть выведен из работы оперативно путем активации сигнала «ОВ КЦТ», при этом осуществляется вывод функций контроля токовых цепей каждой из сторон силового трансформатора (ВН, НН1, НН2).

\item
Ввод функции контроля токовых цепей каждой стороны в работу на этапе параметрирования устройства осуществляется своим программным переключателем SGF1 «Ввод\_функции». Информация о введенной функции отображается выходным сигналом для каждой функции («КЦТ ВН: Ввод», «КЦТ НН1: Ввод», «КЦТ НН2: Ввод»).

\item
Функция контроля токовых цепей любой стороны силового трансформатора (ВН, НН1, НН2) «КЦТ ВН (НН1, НН2)» может быть выведена из работы оперативно путем активации сигнала «ОВ КЦТ ВН», «ОВ КЦТ НН1» или «ОВ КЦТ НН2» соответственно, при этом активизируется выходной сигнал «Оперативный вывод» для соответствующих функций.

\item
Алгоритм работы КЦТ выполнен в соответствии с рисунком \ref{kzt:img1}. В таблице \ref{kzt:tbl1} приведены параметры, необходимые для настройки КЦТ.

\small
\begin{longtable}{|>{\centering\arraybackslash}m{5.3cm}|>{\centering\arraybackslash}m{3.3cm}|>{\centering\arraybackslash}m{4.2cm}|>{\centering\arraybackslash}m{1.8cm}|>{\centering\arraybackslash}m{1cm}|}
\caption{Параметры для настройки КЦТ\hfill\vspace{-0.5\baselineskip}}\label{kzt:tbl1}\\ 
\hline
\rowcolor{gray!30}
Параметр (Параметр на ИЧМ) & Условное обозначение на схеме & Значение/ Диапазон & Единица измерения & Шаг \\ 
\hline
\endfirsthead
%\caption{\emph{Продолжение\hfill\vspace{-0.5\baselineskip}}} \\ % Отступ обозначения таблицы от таблицы и выравнивание надписи слева
\caption*{\hspace{3pt}\emph{Продолжение таблицы \ref{kzt:tbl1}\hfill\vspace{-0.5\baselineskip}}} \\ % сделано по ГОСТ 2.105 п.6.8.7
\hline
\rowcolor{gray!30}
Параметр (Параметр на ИЧМ) & Условное обозначение на схеме & Значение/ Диапазон & Единица измерения & Шаг \\ 
%\hline
\endhead
%\hline
\endfoot
\hline
\endlastfoot
%==+t1*RCTR1|CTR_UIRZ
\centering Ввод функции в работу (Ввод\_функции) & \centering SGF1 & \centering 0 = Не предусмотрено\\1 = Предусмотрено & \centering -- & \centering \arraybackslash -- \\
\hline
\centering Режим контроля обрыва провода (Контр\_обр\_пров) & \centering SGF2 & \centering 0 = Не предусмотрено\\1 = Предусмотрено & \centering -- & \centering \arraybackslash -- \\
\hline
\centering Номинальный ток токового входа терминала (Iном) & \centering -- & \centering 0,2 ... 5,0 & \centering А & \centering \arraybackslash 0,1 \\
\hline
\centering Минимальное значение фазного тока (Iмин) & \centering Iмин< & \centering 0,05 ... 1,00 & \centering о.е. & \centering \arraybackslash 0,01 \\
\hline
\centering Коэффициент симметрии (Kсим) & \centering Iмин/Iмакс< & \centering 0,10 ... 0,95 & \centering о.е. & \centering \arraybackslash 0,01 \\
\hline
\centering Минимальная величина максимального из фазных токов (LIсим) & \centering Iмакс> & \centering 0,01 ... 2,00 & \centering о.е. & \centering \arraybackslash 0,01 \\
\hline
\centering Выдержка времени на срабатывание по критерию обрыва фазы (Tобрыв) & \centering T1 & \centering 0,00 ... 100,00 & \centering с & \centering \arraybackslash 0,01 \\
\hline
\centering Выдержка времени на срабатывание по асимметрии токов (Tасим) & \centering T2 & \centering 0,00 ... 100,00 & \centering с & \centering \arraybackslash 0,01 \\
%===t1
\end{longtable}
\normalsize

\vspace{1mm}
\begin{figure}[H]
\centering

\includegraphics[width=0.9\textwidth,height=0.9\textheight,keepaspectratio]{\fbpath/23. КЦТ ДЗТ/_latex/img/1.pdf}
\captionof{figure}{Функциональная схема логики КЦТ}
\label{kzt:img1}
\end{figure}

%%%%%%%%%%%%%%%%%%%%%%%%%%%%%%%%%%%%%%%%%%%%%%%%%%%%%%%%%%%%%%%%%%%%%% КОНТРОЛЬ ИСПРАВНОСТИ ТОКОВЫХ ЦЕПЕЙ СТОРОНЫ %%%%%%%%%%%%%%%%%%%%%%%%%%%%%%%%%%%%%%%%%%%%%%%%%%

%\emph{Функция контроля исправности токовых цепей стороны трансформатора «КЦТст»}

\item
Функция контроля исправности токовых цепей стороны трансформатора (КЦТст) предназначена для контроля вторичных цепей трансформаторов тока одной из сторон силового трансформатора (возможен контроль токовых цепей до трех сторон). Разомкнутые или короткозамкнутые цепи трансформаторов тока могут вызвать как ложное срабатывание, так и несрабатывание защит.
\item
КЦТст обнаруживает обрывы цепей, короткие замыкания в цепях переменного тока по наличию несимметрии. Наличие несимметрии предполагается, если выполняются следующие условия:

\begin{equation*}
\frac{\left | I_{\textsl{мин.изм.}} \right |}{\left | I_{\textsl{макс.изм.}} \right |}< K_{\textsl{сим}},
\end{equation*}
\begin{equation*}
I_{\textsl{макс.изм.}}> LI_{\textsl{сим}},
\end{equation*}
\begin{eqexpl}[25mm]
\item{$I_{\textsl{мин.изм.}}$, $I_{\textsl{макс.изм.}}$ }минимальный и максимальный измеренные фазные токи соответственно;
\item{$K_{\textsl{сим}}$}коэффициент симметрии;
\item{$LI_{\textsl{сим}}$}минимальная величина максимального из фазных токов.
\end{eqexpl}
В случае выполнения данных условий в течение выдержки времени «Tасим» происходит срабатывание функции.
\item
Характеристика срабатывания КЦТ стороны приведена на рисунке \ref{kitz:char}.
\vspace{1mm}
\begin{figure}[H]
\centering
\includegraphics[width=0.5\textwidth,height=0.5\textheight,keepaspectratio]{\fbpath/23. КЦТ ДЗТ/_latex/img/charkitz.pdf}
\captionof{figure}{Характеристика срабатывания функции контроля токовых цепей стороны трансформатора}\label{kitz:char}
\end{figure}

\end{enumerate}
\FloatBarrier % Все float'ы ДО этой строки будут размещены здесь.